\documentclass[conference]{IEEEtran}
\IEEEoverridecommandlockouts

\usepackage{cite}
\usepackage{amsmath,amssymb,bm}
\usepackage{graphicx}
\usepackage{xcolor}
\usepackage{url}
\usepackage{booktabs}
\usepackage{multirow}

% ---------------------------------------------------------------------------
% Every number quoted in the text comes from generated/results.tex, which is
% written by  python src/make_report_w4.py  from the files in runs/workshop4/.
% Re-running the experiments therefore updates the report automatically.
% \figw4{width}{file}{caption}{label}: full-width figure, only if the PNG exists.
% ---------------------------------------------------------------------------
% Auto-generated by src/make_report_w4.py - do not edit by hand
\newcommand{\zstar}{8}
\newcommand{\betastar}{4}
\newcommand{\epochs}{50}
\newcommand{\daenoise}{0.2}
\newcommand{\aeMSE}{0.0045}
\newcommand{\aeRsq}{0.996}
\newcommand{\aeCorr}{0.998}
\newcommand{\vaeMSE}{0.5106}
\newcommand{\vaeRsq}{0.509}
\newcommand{\vaeCorr}{0.726}
\newcommand{\aeMSEsafe}{0.0044}
\newcommand{\aeMSEfail}{0.0057}
\newcommand{\vaeMSEsafe}{0.5118}
\newcommand{\vaeMSEfail}{0.4885}
\newcommand{\aeTopAnom}{5.6\%}
\newcommand{\vaeTopAnom}{16.7\%}
\newcommand{\aeTopXdot}{0.87}
\newcommand{\vaeTopXdot}{0.88}
\newcommand{\allXdot}{0.47}
\newcommand{\aeTopX}{0.70}
\newcommand{\vaeTopX}{1.57}
\newcommand{\allX}{0.43}
\newcommand{\aeTopFar}{7.4\%}
\newcommand{\vaeTopFar}{68.5\%}
\newcommand{\allFar}{8.8\%}
\newcommand{\aeTopThetadot}{0.42}
\newcommand{\vaeTopThetadot}{0.76}
\newcommand{\allThetadot}{0.35}
\newcommand{\testAnomRate}{5.4\%}
\newcommand{\vaeKL}{1.01}
\newcommand{\vaeActive}{4}
\newcommand{\aeAUC}{0.856}
\newcommand{\aePR}{0.214}
\newcommand{\aeFone}{0.261}
\newcommand{\aeRec}{0.910}
\newcommand{\aePrec}{0.152}
\newcommand{\vaeAUC}{0.920}
\newcommand{\vaePR}{0.372}
\newcommand{\vaeFone}{0.411}
\newcommand{\vaeklAUC}{0.948}
\newcommand{\vaeklPR}{0.503}
\newcommand{\vaeklFone}{0.508}
\newcommand{\daeAUC}{0.898}
\newcommand{\daePR}{0.370}
\newcommand{\supAUC}{0.997}
\newcommand{\supPR}{0.957}
\newcommand{\supFone}{0.848}
\newcommand{\gbAUC}{0.996}
\newcommand{\gbPR}{0.935}
\newcommand{\gbFone}{0.860}
\newcommand{\lrAUC}{0.981}
\newcommand{\lrPR}{0.791}
\newcommand{\aeQrec}{0.114}
\newcommand{\aeQprec}{0.363}
\newcommand{\vaeklQrec}{0.369}
\newcommand{\vaeklQprec}{0.569}
\newcommand{\aeAllAUC}{0.683}
\newcommand{\silRaw}{0.208}
\newcommand{\silAE}{0.145}
\newcommand{\silVAE}{0.234}
\newcommand{\silSup}{0.653}
\newcommand{\probeRaw}{0.981}
\newcommand{\probeAE}{0.928}
\newcommand{\probeVAE}{0.959}
\newcommand{\probeSup}{0.997}
\newcommand{\probePRraw}{0.791}
\newcommand{\probePRAE}{0.436}
\newcommand{\probePRVAE}{0.544}
\newcommand{\probePRSup}{0.960}
\newcommand{\aeNoisyMSE}{0.0337}
\newcommand{\daeNoisyMSE}{0.0154}
\newcommand{\aeCleanMSE}{0.0042}
\newcommand{\daeCleanMSE}{0.0069}
\newcommand{\noiseVar}{0.0400}
\newcommand{\genWvae}{0.348}
\newcommand{\genWae}{0.510}
\newcommand{\genConsVae}{0.032}
\newcommand{\genConsAe}{0.013}
\newcommand{\genBeyond}{0.1\%}
\newcommand{\activeHalf}{8}
\newcommand{\klHalf}{6.62}
\newcommand{\mseHalf}{0.0342}
\newcommand{\silbetaHalf}{0.100}
\newcommand{\activeOne}{7}
\newcommand{\klOne}{4.55}
\newcommand{\mseOne}{0.0755}
\newcommand{\silbetaOne}{0.135}
\newcommand{\activeFour}{4}
\newcommand{\klFour}{1.01}
\newcommand{\mseFour}{0.5106}
\newcommand{\silbetaFour}{0.234}
\newcommand{\trainRowsAll}{24,782}
\newcommand{\trainRowsNormal}{23,402}
\newcommand{\aeMSEzTwo}{0.2377}
\newcommand{\aeAUCzTwo}{0.794}
\newcommand{\aeMSEzFour}{0.0573}
\newcommand{\aeAUCzFour}{0.785}
\newcommand{\aeMSEzEight}{0.0045}
\newcommand{\aeAUCzEight}{0.856}

% Allow up to two full-width figures at the top of a page (many figures in this report)
\renewcommand{\dbltopfraction}{0.9}
\renewcommand{\dblfloatpagefraction}{0.8}
\renewcommand{\topfraction}{0.9}
\renewcommand{\textfraction}{0.1}
\setcounter{dbltopnumber}{2}
\newcommand{\figw}[4]{%
  \IfFileExists{figures/#2.png}{%
  \begin{figure*}[!t]
    \centering
    \includegraphics[width=#1]{figures/#2.png}
    \caption{#3}
    \label{#4}
  \end{figure*}}{}}

\begin{document}

\title{Autoencoders and Variational Autoencoders for Unsupervised Failure Detection in CartPole Trajectories}

\author{\IEEEauthorblockN{Miguel Angel Torres Ramirez}
\IEEEauthorblockA{Dept. of Computer Engineering\\
Universidad Distrital Francisco Jos\'e de Caldas\\
Email: migatorresr@udistrital.edu.co}
\and
\IEEEauthorblockN{Cristian David Romero Gil}
\IEEEauthorblockA{Dept. of Computer Engineering\\
Universidad Distrital Francisco Jos\'e de Caldas\\
Email: cdromerog@udistrital.edu.co}
\and
\IEEEauthorblockN{Omar Esteban Pe\~na Sanchez}
\IEEEauthorblockA{Dept. of Computer Engineering\\
Universidad Distrital Francisco Jos\'e de Caldas\\
Email: oepenas@udistrital.edu.co}}

\maketitle

\begin{abstract}
This report studies what unsupervised representation learning reveals about the CartPole
failure-prediction dataset introduced in Workshop~1. We train fully-connected
autoencoders (AE) with latent sizes $z\in\{2,4,8\}$ and variational autoencoders (VAE)
with $\beta\in\{0.5,1,4\}$ on the same 16 engineered features and the same episode-level
split, without using the failure labels. Pre-failure states are treated as anomalies:
models trained only on safe states score each test state by its reconstruction error.
The deterministic AE reconstructs the data almost perfectly ($R^2=\aeRsq$ at $z=\zstar$) but
also reconstructs pre-failure states, reaching only ROC-AUC $=\aeAUC$. The KL-regularised
VAE is a better detector: the reconstruction-plus-KL score reaches ROC-AUC $=\vaeklAUC$ and
PR-AUC $=\vaeklPR$, without labels. Supervised models remain clearly superior
(Gradient Boosting: ROC-AUC $=\gbAUC$, PR-AUC $=\gbPR$). The latent-space analysis explains
this gap: unsupervised codes organise the data by the dominant physical variables (cart
position and velocity), whereas a supervised bottleneck collapses onto the single direction
that matters for failure (balanced silhouette \silSup{} versus \silVAE{} for the VAE).
\end{abstract}

\begin{IEEEkeywords}
autoencoder, variational autoencoder, $\beta$-VAE, anomaly detection, representation learning, latent space, CartPole.
\end{IEEEkeywords}

%% ======================================================================
\section{Introduction}
Workshop~1 framed failure prediction in CartPole \cite{barto1983,brockman2016} as a
supervised problem: a classifier receives the state at time $t$ and must raise an alarm
during the last five steps before the pole falls or the cart leaves the track. Supervision,
however, has a cost. Failures are rare (5.5\% of the states), they must be
labelled after the fact, and a monitoring system deployed on a new controller may never have
seen one. \emph{Unsupervised} learning offers an alternative: learn what normal operation
looks like and flag whatever does not fit.

Autoencoders learn a compressed code by reconstructing their input
\cite{hinton2006,goodfellow2016}; variational autoencoders add a probabilistic latent space
regularised towards a prior \cite{kingma2014,doersch2016}. Both give two tools for this
domain: a latent space to inspect, and a reconstruction error that can serve as an anomaly
score \cite{chalapathy2019}. The central question of this workshop is: \emph{what structure
do unsupervised models reveal about the dataset, and how can these representations enhance
our understanding and the downstream task?}

The contributions are:
(i) an AE and a $\beta$-VAE adapted to tabular control data, with a latent-size sweep and a
$\beta$ sweep;
(ii) a reconstruction-based failure detector trained only on safe states, compared with the
Workshop~1 baselines and a supervised neural reference;
(iii) a latent-space analysis (t-SNE, UMAP, silhouette, linear probes and per-dimension
correlations) contrasting unsupervised and supervised representations; and
(iv) two extensions: a denoising AE and an evaluation of the VAE as a generative model.

%% ======================================================================
\section{Data and Problem Setup}

\subsection{Dataset and Inputs}
We reuse the Workshop~1 dataset without changes: 35,874 states from 400 CartPole-v1
episodes, generated with a mixed noisy-expert/random policy. Each input is the vector of the
16 engineered features of Workshop~1. These are the raw state $(x,\dot{x},\theta,\dot\theta)$
and the action, the distances to the failure limits, first differences, 3-step rolling means
and time-to-limit estimates. None of them uses future information. Features are z-scored
with statistics from the training episodes only. The episode-level 70/15/15 split stored in
\texttt{data/splits.json} is reused, so every result is directly comparable with
Workshop~1 and there is no leakage between consecutive states.

\subsection{Normal and Anomalous States}
The dataset has a natural normal/abnormal partition: a state is \emph{normal} if
\texttt{label}$=0$ (safe) and \emph{anomalous} if \texttt{label}$=1$ (failure within five
steps). The test set contains 5,330 states, \testAnomRate{} of them anomalous. Each
unsupervised model is trained in two regimes:
\begin{itemize}
\item \textbf{All rows} (\trainRowsAll{} training states). These models are used to measure
reconstruction quality and to analyse the latent space.
\item \textbf{Normal rows only} (\trainRowsNormal{} training states). These models are used
for anomaly detection, as the workshop prescribes.
\end{itemize}
Labels are used only to evaluate the models, to tune thresholds on validation and to select
hyper-parameters on validation. They are never used to train an unsupervised model.

%% ======================================================================
\section{Methods}

\subsection{Autoencoder}
Because the data are tabular, encoder and decoder are fully-connected networks
(Table~\ref{tab:arch}). The decoder mirrors the encoder and ends in a \emph{linear} layer,
since the targets are unbounded z-scores and a Tanh or Sigmoid would clip them. The loss is
the mean squared error (MSE). The workshop suggests $z\in\{32,64,128\}$ and a latent size of
about 5--10\% of the input dimension. With 16 inputs, the first suggestion would give an
\emph{over-complete} code that can learn the identity, while 5--10\% means 1--2 units. We
therefore sweep $z\in\{2,4,8\}$, which brackets the guideline and stays below the input
dimension.

\begin{table}[t]
\caption{Network Architectures (ReLU Hidden Activations)}
\label{tab:arch}
\centering
\begin{tabular}{ll}
\hline
\textbf{Block} & \textbf{Layers} \\ \hline
AE encoder & $16 \rightarrow 64 \rightarrow 32 \rightarrow z$ \\
VAE encoder & $16 \rightarrow 64 \rightarrow 32 \rightarrow (\bm\mu, \log\bm\sigma^2)\in\mathbb{R}^{z}\times\mathbb{R}^{z}$ \\
Decoder (both) & $z \rightarrow 32 \rightarrow 64 \rightarrow 16$ (linear output) \\
Supervised MLP & $16 \rightarrow 64 \rightarrow 32 \rightarrow z$ (ReLU) $\rightarrow 1$ (logit) \\
\hline
\end{tabular}
\end{table}

\subsection{Variational Autoencoder}
The VAE encoder outputs the parameters of a diagonal Gaussian posterior
$q_\phi(\mathbf z|\mathbf x)=\mathcal N(\bm\mu,\mathrm{diag}\,\bm\sigma^2)$. A latent sample
is obtained with the \emph{reparameterisation trick}
\begin{equation}
\mathbf z=\bm\mu+\bm\sigma\odot\bm\epsilon,\qquad \bm\epsilon\sim\mathcal N(\mathbf 0,\mathbf I),
\end{equation}
which moves the randomness to an external input. The sample then becomes a deterministic,
differentiable function of $\bm\mu$ and $\bm\sigma$, so gradients flow through the encoder.
The network minimises the negative evidence lower bound (ELBO) with a weight $\beta$ on the
KL term \cite{higgins2017}:
\begin{align}
\mathcal L &= \underbrace{\textstyle\sum_{j=1}^{16}(x_j-\hat x_j)^2}_{\mathcal L_{\text{recon}}}
 + \beta\,\mathcal L_{\text{KL}}, \\
\mathcal L_{\text{KL}} &= -\tfrac12\textstyle\sum_{k=1}^{z}\left(1+\log\sigma_k^2-\mu_k^2-\sigma_k^2\right).
\end{align}
The reconstruction term is the squared error \emph{summed} over features, which corresponds
to a Gaussian decoder with fixed variance. This keeps it on the same per-sample scale as the
KL term, so $\beta$ has a clear meaning. To prevent posterior collapse (KL vanishing), we use
KL annealing \cite{bowman2016}: $\beta$ rises linearly from 0 to its target over the first
10 epochs. At evaluation time the VAE encodes and reconstructs with the posterior mean
$\bm\mu$. Collapse is diagnosed with the number of \emph{active units}: dimensions whose
$\mathrm{Var}_{\mathbf x}[\mu_k]$ exceeds 0.01 \cite{burda2016}.

\subsection{Denoising Autoencoder (Extension)}
Following \cite{vincent2010}, a denoising AE with the selected latent size receives
$\tilde{\mathbf x}=\mathbf x+\mathcal N(0,\daenoise^2\mathbf I)$ and is trained to reconstruct
the clean $\mathbf x$. The noise acts as a regulariser against the identity mapping.

\subsection{Supervised Reference Model}
The workshop compares the unsupervised codes with the representation of the Workshop~2
supervised network. That model is not part of the repository, so we train a reference
supervised MLP with the same trunk and a bottleneck of the same width $z$
(Table~\ref{tab:arch}). It uses binary cross-entropy weighted by the class ratio and the same
data and split. Its frozen bottleneck is the \emph{supervised} representation; results from
it are marked with $\dagger$.

\subsection{Training and Model Selection}
All models use Adam \cite{kingma2015} (learning rate $10^{-3}$, batch 64) for \epochs{}
epochs with seed 42. We log train and validation losses every epoch and keep the best
validation epoch. For the VAE, the best epoch is chosen with the target-$\beta$ objective and
only after the warm-up. The supervised MLP keeps the epoch with the best validation PR-AUC.

Selection uses validation data only. First, $z^\star$ is the latent size whose
normal-rows AE has the highest validation anomaly ROC-AUC. Second, the VAE uses $z^\star$,
and $\beta^\star$ is selected the same way. This gives $z^\star=\zstar$ and
$\beta^\star=\betastar$. The \emph{best} checkpoints (\texttt{ae\_best.pt},
\texttt{vae\_best.pt}) are the all-rows models with these settings.

\subsection{Evaluation Protocol}
\textbf{Reconstruction.} On the test set we report MSE in standardised units, the explained
variance $R^2$ and the mean Pearson correlation $\bar\rho$ between each feature and its
reconstruction. These are the continuous-data metrics the workshop suggests; SSIM and PSNR
do not apply to tabular data.

\textbf{Anomaly detection.} The score is the per-state reconstruction MSE. For the VAE we
also use $\mathcal L_{\text{recon}}+\mathcal L_{\text{KL}}$. We report the threshold-free
ROC-AUC and PR-AUC. Precision, recall and F1 are reported at the threshold that maximises F1
on \emph{validation}, the same rule as in Workshop~1. As a label-free alternative, we also
use the 99th percentile of the validation normal scores.

\textbf{Latent space.} Test states are encoded (with $\bm\mu$ for the VAE) and compared with
the input features and the supervised bottleneck. We use four measurements:
(a) the silhouette score \cite{rousseeuw1987} of the label partition, after z-scoring each
dimension. It is computed on all rows and, because of the 17:1 imbalance, on ten
class-balanced subsamples.
(b) A \emph{linear probe} \cite{alain2016}: a logistic regression fitted on training codes
and evaluated on test, which measures how much failure information is linearly available.
(c) 2-D t-SNE \cite{vandermaaten2008} and UMAP \cite{mcinnes2018} maps.
(d) The correlation between each latent dimension and the physical variables.

%% ======================================================================
\section{Results}

\figw{\textwidth}{fig-w4-training-curves}{Training curves. (a) AE reconstruction MSE for each
latent size. (b) VAE reconstruction MSE and (c) validation KL divergence for each $\beta$
($z=z^\star$); the shaded band is the KL warm-up.}{fig:curves}

\subsection{Training Dynamics}
Fig.~\ref{fig:curves} shows stable training without overfitting: the train and validation
curves stay close in every run. The AE error falls sharply with the latent size (test MSE \aeMSEzTwo, \aeMSEzFour{} and
\aeMSEzEight{} for $z=2,4,8$). With $z=8$ the code is only half the input width, yet the
reconstruction is nearly perfect, because the 16 engineered features are highly redundant. Several of them are deterministic functions of
the raw state, for example $|\theta|$ and $0.2095-|\theta|$.

The VAE curves show the annealing mechanism at work. While $\beta$ is still small, the
reconstruction error drops quickly. When the KL weight reaches its target, the error rises
until the model settles on a trade-off whose position depends on $\beta$. For $\beta=4$ the
KL term falls to \klFour{} nats per sample.

\subsection{Reconstruction Quality}
Table~\ref{tab:recon} summarises the test reconstruction. The selected AE
($z=\zstar$) reaches MSE $=\aeMSE$ and $\bar\rho=\aeCorr$. The selected VAE
($\beta=\betastar$) is far less accurate (MSE $=\vaeMSE$, $R^2=\vaeRsq$). This is the
expected rate-distortion trade-off: as $\beta$ grows, the model pays more for each nat of
information in $\mathbf z$ and keeps only the dominant factors. Only \vaeActive{} of its
\zstar{} latent dimensions remain active (Fig.~\ref{fig:beta}b), so the posterior has
partially collapsed. The remaining dimensions encode the variables with the most variance.

\begin{table}[t]
\caption{Test Reconstruction Quality (All-Rows Models). MSE$_0$/MSE$_1$: Safe/Failure States; AU: Active Units; $\star$: Selected}
\label{tab:recon}
\centering
\setlength{\tabcolsep}{3.2pt}
\resizebox{\columnwidth}{!}{\begin{tabular}{lccccccc}
\hline
\textbf{Model} & \textbf{MSE} & $\mathbf{R^2}$ & $\bar{\boldsymbol{\rho}}$ & \textbf{MSE$_0$} & \textbf{MSE$_1$} & \textbf{KL} & \textbf{AU} \\ \hline
AE, $z=2$ & 0.2377 & 0.771 & 0.874 & 0.2317 & 0.3418 & -- & -- \\
AE, $z=4$ & 0.0573 & 0.945 & 0.972 & 0.0570 & 0.0635 & -- & -- \\
AE, $z=8$$^\star$ & \textbf{0.0045} & \textbf{0.996} & \textbf{0.998} & 0.0044 & 0.0057 & -- & -- \\
\hline
VAE, $z=8$, $\beta=0.5$ & 0.0342 & 0.967 & 0.985 & 0.0331 & 0.0531 & 6.62 & 8/8 \\
VAE, $z=8$, $\beta=1$ & 0.0755 & 0.927 & 0.968 & 0.0733 & 0.1122 & 4.55 & 7/8 \\
VAE, $z=8$, $\beta=4$$^\star$ & 0.5106 & 0.509 & 0.726 & 0.5118 & 0.4885 & 1.01 & 4/8 \\
\hline
\end{tabular}
}
\end{table}

\figw{\textwidth}{fig-w4-reconstructions}{Sample reconstructions of eight test states (four safe,
four within five steps of failure). Columns are the 16 standardised features. From left to
right: original, AE reconstruction ($z=z^\star$), VAE reconstruction ($\beta=\beta^\star$)
and element-wise absolute-error maps on a common scale.}{fig:samples}

Fig.~\ref{fig:samples} makes the difference visible. The AE copies every feature, including
the binary action. The VAE keeps only the coarse pattern of each state, such as the sign of $\theta$ and
$\dot\theta$ and whether the pole is close to its limit. It smooths away the binary action,
the change in angular velocity and the time-to-limit estimates, which are the features with
the largest errors in Fig.~\ref{fig:errors}a.

\figw{\textwidth}{fig-w4-error-analysis}{Error analysis on the test set. (a) MSE per feature.
(b) Median reconstruction error versus the number of steps left before a failure, for the
AE trained on all rows and on normal rows only. (c) Distribution of the per-state error by
class.}{fig:errors}

\textbf{Which states are hardest to reconstruct?} Errors grow as the system approaches
failure (Fig.~\ref{fig:errors}b). The increase is steeper for the model that never saw a
pre-failure state, which is the signal used for anomaly detection in the next subsection.
The hardest states are not necessarily pre-failure states. The 1\% of test states with the
largest AE error are fast-moving carts: their mean $|\dot x|$ is \aeTopXdot{}~m/s, compared
with \allXdot{}~m/s overall. Only \aeTopAnom{} of them are pre-failure states, close to the
\testAnomRate{} base rate. For the VAE, the hardest states are carts far from the centre:
\vaeTopFar{} have $|x|>1.5$~m, compared with \allFar{} overall, and \vaeTopAnom{} are
pre-failure states. In both cases the errors concentrate in sparsely populated regions of the
state space, whether or not those regions lead to a failure. The AE trained on all rows
reconstructs failure states almost as well as safe ones (Table~\ref{tab:recon}, MSE$_1$
versus MSE$_0$). A flexible AE generalises to anything that
lies on the same manifold, so a good reconstruction does not by itself make a good detector.

\subsection{Anomaly Detection}
Table~\ref{tab:anomaly} is the cumulative comparison requested by the workshop: the
Workshop~1 baselines, the supervised reference network and the Workshop~4 detectors on the
same test episodes. Table~\ref{tab:sweep} reports the sweep over latent size, $\beta$ and
training regime. Fig.~\ref{fig:anomaly} shows the curves.

\begin{table}[t]
\caption{Cumulative Test Results for Failure / Anomaly Detection (Threshold Selected on Validation). Bold: Best Unsupervised; $\dagger$: Supervised Reference Re-Trained in This Workshop}
\label{tab:anomaly}
\centering
\setlength{\tabcolsep}{4pt}
\begin{tabular}{llccccc}
\hline
& \textbf{Model} & \textbf{ROC} & \textbf{PR} & \textbf{P} & \textbf{R} & \textbf{F1} \\ \hline
W1 & Logistic Reg. & 0.981 & 0.791 & 0.669 & 0.703 & 0.686 \\
W1 & Random Forest & 0.994 & 0.923 & 0.812 & 0.852 & 0.832 \\
W1 & Gradient Boost. & 0.996 & 0.935 & 0.834 & 0.886 & 0.860 \\
W2$^\dagger$ & MLP (W2 ref.) & 0.997 & 0.957 & 0.770 & 0.945 & 0.848 \\
\hline
W4 & AE & 0.856 & 0.214 & 0.152 & \textbf{0.910} & 0.261 \\
W4 & VAE: recon. & 0.920 & 0.372 & 0.375 & 0.455 & 0.411 \\
W4 & VAE: recon.\,+\,KL & \textbf{0.948} & \textbf{0.503} & \textbf{0.497} & 0.521 & \textbf{0.508} \\
W4 & Denoising AE & 0.898 & 0.370 & 0.412 & 0.469 & 0.439 \\
\hline
\end{tabular}

\end{table}

\begin{table}[t]
\caption{Anomaly ROC-AUC by Latent Size, $\beta$ and Training Regime (Test in Parentheses: VAE Recon.\,+\,KL Score)}
\label{tab:sweep}
\centering
\setlength{\tabcolsep}{5pt}
\begin{tabular}{lccc}
\hline
& \multicolumn{2}{c}{\textbf{Normal rows only}} & \textbf{All rows} \\
\textbf{Model} & \textbf{Val} & \textbf{Test} & \textbf{Test} \\ \hline
AE, $z=2$ & 0.793 & 0.794 & 0.650 \\
AE, $z=4$ & 0.785 & 0.785 & 0.572 \\
AE, $z=8$$^\star$ & 0.825 & 0.856 & 0.683 \\
\hline
VAE, $\beta=0.5$ & 0.737 & 0.775 (0.964) & 0.546 \\
VAE, $\beta=1$ & 0.798 & 0.813 (0.954) & 0.555 \\
VAE, $\beta=4$$^\star$ & 0.898 & 0.920 (0.948) & 0.469 \\
\hline
\end{tabular}

\end{table}

\figw{\textwidth}{fig-w4-anomaly}{Anomaly detection on the test set. (a) ROC and (b)
precision-recall curves of the unsupervised detectors (trained on safe states only) and
the supervised reference; the grey line in (b) is the prevalence. (c) Distribution of the AE
anomaly score by class.}{fig:anomaly}

Four findings stand out.
\begin{enumerate}
\item \textbf{Training on normal rows is essential.} The same AE trained on all rows reaches
only ROC-AUC $=\aeAllAUC$, compared with \aeAUC{} when trained on safe states only
(Table~\ref{tab:sweep}). Once pre-failure states appear in the training data, the model
learns to reconstruct them.
\item \textbf{Regularisation beats reconstruction accuracy.} The VAE reconstructs worse
than the AE but detects better: ROC-AUC \vaeAUC{} versus \aeAUC. A strongly regularised
model can only represent the bulk of the normal data, so unusual states are pushed towards
typical ones and leave a larger residual. The denoising AE, also regularised, falls in
between (\daeAUC).
\item \textbf{The KL term is informative.} Adding the KL divergence to the VAE score raises
ROC-AUC to \vaeklAUC{} and PR-AUC to \vaeklPR. Pre-failure states are encoded far from the
prior, so the posterior itself signals novelty, consistent with \cite{schlegl2017}.
\item \textbf{Supervision still wins by a wide margin.} The best unsupervised detector ranks
states well, but its PR-AUC (\vaeklPR) is far below Gradient Boosting (\gbPR) and the
supervised MLP (\supPR). Its F1 at the validation threshold is \vaeklFone{}, versus \gbFone{}.
The score distributions of the two classes overlap (Fig.~\ref{fig:anomaly}c): many safe
states (e.g.\ large but recoverable angles under the random policy) look as unusual as the
last steps before a failure. The label-free 99th-percentile threshold confirms the
trade-off: the AE then has precision \aeQprec{} and recall \aeQrec.
\end{enumerate}

\subsection{Latent Space Structure}

\figw{\textwidth}{fig-w4-latent-tsne}{t-SNE maps of the test states in the input space, the AE code,
the VAE posterior mean and the supervised bottleneck. Top: coloured by label. Bottom:
coloured by $|\theta|$.}{fig:tsne}

\figw{\textwidth}{fig-w4-latent-umap}{UMAP maps of the same representations as Fig.~\ref{fig:tsne}.}{fig:umap}

\begin{table}[t]
\caption{Separability of the Failure Class in Each Representation (Test). Silhouette on the Full Set, on Balanced Subsamples and on the 2-D Maps; Linear Probe ROC-AUC and PR-AUC}
\label{tab:latent}
\centering
\setlength{\tabcolsep}{3pt}
\resizebox{\columnwidth}{!}{\begin{tabular}{lccccccc}
\hline
& & \multicolumn{4}{c}{\textbf{Silhouette}} & \multicolumn{2}{c}{\textbf{Linear probe}} \\
\textbf{Space} & \textbf{Dim} & \textbf{All} & \textbf{Bal.} & \textbf{t-SNE} & \textbf{UMAP} & \textbf{ROC} & \textbf{PR} \\ \hline
Input features & 16 & 0.336 & 0.208 & -0.118 & -- & 0.981 & 0.791 \\
AE code & 8 & 0.317 & 0.145 & -0.015 & -- & 0.928 & 0.436 \\
VAE $\mu$ & 8 & 0.370 & 0.234 & 0.030 & -- & 0.959 & 0.544 \\
MLP (W2 ref.)$^\dagger$ & 8 & 0.737 & 0.653 & 0.095 & -- & 0.997 & 0.960 \\
\hline
\end{tabular}
}
\end{table}

\figw{0.92\textwidth}{fig-w4-latent-correlation}{Absolute Pearson correlation between each latent
dimension and interpretable variables (test set).}{fig:corr}

\figw{0.85\textwidth}{fig-w4-beta}{Effect of $\beta$ on the VAE ($z=z^\star$, all rows).
(a) Reconstruction MSE versus KL (rate-distortion). (b) KL per latent dimension; values close
to zero indicate collapsed dimensions. (c) Balanced silhouette of the failure class and
anomaly ROC-AUC of the normal-rows model.}{fig:beta}

\textbf{What do the autoencoders discover without labels?} The t-SNE maps
(Fig.~\ref{fig:tsne}) of the input space and of the AE code are fragmented into many small
clusters. These are trajectory segments with similar cart position and velocity, and the
failure states form several separate islands. The colouring by $|\theta|$ shows that these
islands are the high-angle regions of the state space. The correlation analysis
(Fig.~\ref{fig:corr}) makes this concrete. Each AE and VAE dimension is a mixture of $x$,
$\dot x$, $\theta$ and $\dot\theta$, and their correlation with the label is low. The
autoencoders find the \emph{physical} structure of the system, organised by the variables
that carry the most variance, which are not the variables that decide failure.

The VAE map is somewhat smoother, and its balanced silhouette (\silVAE) is higher than that of the AE (\silAE) and of the input space
(\silRaw). The $\beta$ study (Fig.~\ref{fig:beta}) explains why: increasing $\beta$ removes
the low-variance details (fewer active units: \activeHalf, \activeOne{} and \activeFour{} for
$\beta=0.5,1,4$) and raises both the class separability (silhouette \silbetaHalf{},
\silbetaOne{} and \silbetaFour) and the anomaly AUC. On this dataset, the more compressed and
disentangled code \cite{higgins2017} is also the more useful one.

\subsection{Supervised versus Unsupervised Representations}
The supervised bottleneck is far more organised. Its t-SNE map is a single one-dimensional
curve ordered by risk, with the failure states at one end (Fig.~\ref{fig:tsne}). Its
balanced silhouette is \silSup{} (Table~\ref{tab:latent}). Almost every one of its
dimensions correlates with the label and with $|\theta|$, and none with the cart position
(Fig.~\ref{fig:corr}). In effect, the labels told the network which single direction of the
state space matters and allowed it to discard everything else.

The linear probes quantify the value of each representation for the downstream task. A
logistic regression reaches ROC-AUC \probeSup{} on the supervised code, \probeRaw{} on the
raw features (the Workshop~1 logistic baseline), \probeVAE{} on the VAE code and
\probeAE{} on the AE code. Compression without labels therefore \emph{loses} part of the
failure information. The reconstruction objective treats all variance as equally important,
whereas the failure signal is concentrated in a few directions (mainly the angle and its rate
near the limit) that carry only a small share of the total variance. The lesson about label supervision is that it is not only a training signal
but a \emph{relevance} signal. Unsupervised models answer ``what varies?'', whereas
supervised ones answer ``what matters?''. In this domain the two questions have different
answers.

\subsection{Extensions: Denoising and Generation}
\textbf{Denoising.} With inputs corrupted by Gaussian noise ($\sigma=\daenoise$), the
denoising AE reconstructs the clean test states with MSE \daeNoisyMSE{}, compared with
\aeNoisyMSE{} for the plain AE and \noiseVar{} for doing nothing. On clean inputs the price is
small (\daeCleanMSE{} versus \aeCleanMSE). The denoising objective also improves anomaly
detection (Table~\ref{tab:anomaly}), because it discourages the identity mapping.

\figw{\textwidth}{fig-w4-vae-samples}{Generative properties. Distribution of the state variables in the
test set and in 5,000 samples decoded from $\mathbf z\sim\mathcal N(\mathbf 0,\mathbf I)$ by
the VAE and, as a contrast, by the AE (physical units). Right: consistency between the
generated angle and the generated $|\theta|$ feature.}{fig:gen}

\textbf{Generation.} Sampling $\mathbf z\sim\mathcal N(\mathbf 0,\mathbf I)$ and decoding
(Fig.~\ref{fig:gen}) yields states centred on the real distribution, but too concentrated
around the upright equilibrium. The mean 1-Wasserstein distance over the four state
variables, divided by their standard deviation, is \genWvae{} for the VAE and \genWae{} for
the AE. The generated features are only approximately consistent with each other: the mean
gap between the generated $|\theta|$ feature and the absolute value of the generated angle is
\genConsVae{}~rad (\genConsAe{}~rad for the AE). Because the selected VAE is strongly regularised, it trades sample
diversity for a smooth latent space. Using it for data augmentation would require a smaller
$\beta$ or a decoder variance learned per feature.

%% ======================================================================
\section{Discussion and Limitations}
\begin{itemize}
\item \textbf{Low-dimensional tabular data.} With only 16 partly redundant inputs, an
over-complete or even moderately sized code quickly approaches the identity map. This is
why small latent sizes and strong regularisation work best here, unlike in image domains.
\item \textbf{Single seed and selection variance.} All results come from one seed and one
split. Several validation AUC differences in Table~\ref{tab:sweep} are a few hundredths, so
the selected $z^\star$ and $\beta^\star$ should be read as indicative rather than definitive.
\item \textbf{Supervised reference.} The supervised comparison uses a re-trained MLP with
the same trunk as the autoencoders, not the original Workshop~2 network. The conclusions
concern supervised versus unsupervised training of the same architecture.
\item \textbf{Static inputs.} Each state is encoded independently. A recurrent or windowed
autoencoder over the last few steps could capture the dynamics of an approaching failure
better.
\end{itemize}

%% ======================================================================
\section{Conclusions}
Autoencoders learn a faithful, low-dimensional description of CartPole states, but that
description is organised by the variables with the most variance (cart position and
velocity), not by the proximity to failure. Trained only on safe states, the models turn
reconstruction error into a label-free failure detector. Regularisation is the key
ingredient: the $\beta$-VAE with the reconstruction-plus-KL score reaches ROC-AUC \vaeklAUC,
while the more accurate AE reaches \aeAUC. Supervised models remain far better
(PR-AUC \gbPR{} versus \vaeklPR), and the latent analysis shows why: labels select the one
direction of the state space that matters. In practice, unsupervised models are useful for
understanding the data and for monitoring controllers where failures have never been
observed. When labels exist, the supervised models of Workshops~1 and~2 should be preferred.

\section*{Reproducibility}
Code: \texttt{src/autoencoder.py} (models, losses, training loop), \texttt{src/train\_ae.py}
(experiment plan and logging), \texttt{src/evaluate\_ae.py} (reconstruction, anomaly
detection, denoising, generation), \texttt{src/analyze\_latent.py} (latent analysis) and
\texttt{src/make\_report\_w4.py} (tables and numbers of this report). Checkpoints are in
\texttt{checkpoints/} (\texttt{ae\_best.pt}, \texttt{vae\_best.pt} and the normal-rows and
reference models) and logs in \texttt{runs/workshop4/}. Seed 42; split \texttt{data/splits.json}.

\begin{thebibliography}{99}
\bibitem{barto1983} A.~G. Barto, R.~S. Sutton, and C.~W. Anderson, ``Neuronlike adaptive elements that can solve difficult learning control problems,'' \emph{IEEE Trans. Syst., Man, Cybern.}, vol.~13, no.~5, pp.~834--846, 1983.
\bibitem{brockman2016} G.~Brockman \emph{et al.}, ``OpenAI Gym,'' arXiv:1606.01540, 2016.
\bibitem{hinton2006} G.~E. Hinton and R.~R. Salakhutdinov, ``Reducing the dimensionality of data with neural networks,'' \emph{Science}, vol.~313, no.~5786, pp.~504--507, 2006.
\bibitem{vincent2010} P.~Vincent, H.~Larochelle, I.~Lajoie, Y.~Bengio, and P.-A. Manzagol, ``Stacked denoising autoencoders: Learning useful representations in a deep network with a local denoising criterion,'' \emph{J. Mach. Learn. Res.}, vol.~11, pp.~3371--3408, 2010.
\bibitem{goodfellow2016} I.~Goodfellow, Y.~Bengio, and A.~Courville, \emph{Deep Learning}. MIT Press, 2016, ch.~14.
\bibitem{kingma2014} D.~P. Kingma and M.~Welling, ``Auto-encoding variational Bayes,'' in \emph{Proc. ICLR}, 2014.
\bibitem{doersch2016} C.~Doersch, ``Tutorial on variational autoencoders,'' arXiv:1606.05908, 2016.
\bibitem{higgins2017} I.~Higgins \emph{et al.}, ``$\beta$-VAE: Learning basic visual concepts with a constrained variational framework,'' in \emph{Proc. ICLR}, 2017.
\bibitem{bowman2016} S.~R. Bowman \emph{et al.}, ``Generating sentences from a continuous space,'' in \emph{Proc. CoNLL}, 2016, pp.~10--21.
\bibitem{burda2016} Y.~Burda, R.~Grosse, and R.~Salakhutdinov, ``Importance weighted autoencoders,'' in \emph{Proc. ICLR}, 2016.
\bibitem{chalapathy2019} R.~Chalapathy and S.~Chawla, ``Deep learning for anomaly detection: A survey,'' arXiv:1901.03407, 2019.
\bibitem{schlegl2017} T.~Schlegl, P.~Seeb\"ock, S.~M. Waldstein, U.~Schmidt-Erfurth, and G.~Langs, ``Unsupervised anomaly detection with generative adversarial networks to guide marker discovery,'' in \emph{Proc. IPMI}, 2017, pp.~146--157.
\bibitem{kingma2015} D.~P. Kingma and J.~Ba, ``Adam: A method for stochastic optimization,'' in \emph{Proc. ICLR}, 2015.
\bibitem{rousseeuw1987} P.~J. Rousseeuw, ``Silhouettes: A graphical aid to the interpretation and validation of cluster analysis,'' \emph{J. Comput. Appl. Math.}, vol.~20, pp.~53--65, 1987.
\bibitem{alain2016} G.~Alain and Y.~Bengio, ``Understanding intermediate layers using linear classifier probes,'' arXiv:1610.01644, 2016.
\bibitem{vandermaaten2008} L.~van der Maaten and G.~Hinton, ``Visualizing data using t-SNE,'' \emph{J. Mach. Learn. Res.}, vol.~9, pp.~2579--2605, 2008.
\bibitem{mcinnes2018} L.~McInnes, J.~Healy, and J.~Melville, ``UMAP: Uniform manifold approximation and projection for dimension reduction,'' arXiv:1802.03426, 2018.
\end{thebibliography}

\end{document}
